\section{Function Calls}

\begin{itemize}
  \item Function calls use the VIV-32 link register.
  \item Call instructions write the return address to \(lr\), corresponding to \(r15\).
  \item The callee returns by jumping to the address held in \(lr\), unless it has restored a saved return address first.
\end{itemize}

\section{Argument Passing}

\begin{itemize}
  \item Arguments are passed left-to-right.
  \item The first four word-sized arguments are passed in \(a0\)--\(a3\), corresponding to \(r3\)--\(r6\).
  \item Additional arguments are passed on the stack.
  \item Stack-passed arguments are placed in ascending argument order.
\end{itemize}

\section{Return Values}

\begin{itemize}
  \item Return values up to 32 bits are returned in \(rv0\), corresponding to \(r1\).
  \item Return values up to 64 bits are returned in \(rv0:rv1\).
  \item For 64-bit return values:
  \begin{itemize}
    \item \(rv0\) holds the low 32 bits;
    \item \(rv1\) holds the high 32 bits.
  \end{itemize}
  \item Larger return values are returned indirectly through caller-provided memory.
  \item For return values larger than 64 bits, the caller provides a return buffer and passes its address as a hidden first argument.
\end{itemize}

\section{Aggregate Values}

\begin{itemize}
  \item Small aggregates that fit in one or two words may be passed or returned in registers.
  \item Larger aggregates are passed by address.
  \item Exact aggregate layout rules are deferred until the C ABI is specified.
\end{itemize}

\section{Link Register Preservation}

\begin{itemize}
  \item The link register is overwritten by call instructions.
  \item A leaf function that does not call another function may return directly using \(lr\).
  \item A non-leaf function must save its incoming \(lr\) value if it needs to return to its own caller.
\end{itemize}

\section{Variadic Functions}

\begin{itemize}
  \item Variadic functions use a runtime-defined argument access structure.
  \item Register arguments may be spilled into a register save area in the callee's stack frame.
  \item Additional variadic arguments are read from the stack.
  \item Exact \texttt{va\_list} representation is deferred to the runtime/library ABI.
\end{itemize}

\section{System Call ABI}

\begin{itemize}
  \item System calls are performed using the software trap instruction.
  \item Register \(rv0\), corresponding to \(r1\), contains the system call number.
  \item Registers \(a0\)--\(a3\), corresponding to \(r3\)--\(r6\), contain the first four system call arguments.
  \item Additional system call arguments, if required, are passed by pointer to a memory structure.
  \item The pointer is passed in one of the argument registers.
  \item Register \(rv0\), corresponding to \(r1\), contains the primary system call return value.
  \item Register \(rv1\), corresponding to \(r2\), may contain a secondary return value.
  \item Failed system calls return a negative error code in \(rv0\).
  \item The system call handler may also set the \(E\) flag to indicate failure.
  \item Unless otherwise documented for a specific system call, syscall handlers may clobber caller-saved registers and must preserve callee-saved registers.
  \item VIV-32 v0.1 does not define a separate user mode.
  \item System calls are therefore an operating-environment convention rather than a hardware-enforced privilege transition.
  \item A system call handler returns using the interrupt/trap return mechanism.
  \item System call return restores execution at the instruction following the trap.
  \item The mapping between system call numbers and services is defined by the operating environment, not by the base ISA.
\end{itemize}
